% Continuation of the pruning-depth material. Requires: amsmath, amsthm, booktabs.
% \newtheorem{theorem}{Theorem}
% \newtheorem{lemma}{Lemma}
% \newtheorem{proposition}{Proposition}
% \newtheorem{corollary}{Corollary}
% \newtheorem{definition}{Definition}
%
% Notation assumed from the earlier sections: a block $B$ delegated to a composite
% resource $CR$ of pod $A$; $\mathit{anchors}_V(B)$ the residual (variable) anchors;
% $B(A)$ the delegated part; $\gamma_c(B)$ the connected cover number, the least depth at
% which an operator of the default plan may prune.

\section{The Scope of Delegation}

The pruning analysis gives a criterion for which blocks are worth delegating at all, and
it is sharper than a restriction on shape. We state it, show that it excludes long paths,
and show that long paths can nevertheless be delegated by decomposing them into blocks
that satisfy it.

\begin{definition}[Reducible block]
A block $B$ is \emph{reducible} if $\gamma_c(B) = 1$, that is, if some single pattern of
$B$ binds every residual anchor of $B$:
\[
\exists\, tp \in B \;:\; \mathit{anchors}_V(B) \subseteq \mathit{vars}(tp).
\]
\end{definition}

Reducibility is a property of the block, not of a shape vocabulary. It admits stars of
any size, whose anchor set is a singleton occurring in every pattern; paths of length
two, whose anchors $\{?s, ?o_1\}$ are the variables of $tp_1$; and any block, of whatever
shape, whose anchors happen to be collected in one pattern. It excludes paths of length
three or more, whose anchors are distributed across patterns.

\begin{theorem}[Characterisation]
The delegated part $B(A)$ contributes nothing to the join state of the default plan if
and only if $B$ is reducible.
\end{theorem}

\begin{proof}
By the Generalised Pruning lemma an operator may prune exactly when its mappings bind
every anchor of $B$, and by definition of $\gamma_c$ the shallowest such operator has
depth $\gamma_c(B)$. If $\gamma_c(B) = 1$ that operator scans a single pattern, so the
test is decided on each tuple as it is read, before it is indexed: no mapping of $B(A)$
is ever inserted into the join state. If $\gamma_c(B) > 1$ no depth-$1$ operator can
decide the test, so every tuple of $B(A)$ is indexed, and the mappings composing $B(A)$
are built by at least $\gamma_c(B) - 1$ joins before any of them can be discarded.
\end{proof}

The two cases differ in what pruning removes. At depth $1$ a discarded tuple never enters
the index, never probes, and never produces the intermediate results that would have
descended from it, so the saving is the whole sub-derivation rooted at that tuple and
scales with its fanout. At depth $\gamma_c(B) > 1$ the mapping has already been built,
and discarding it removes one join and the routing that would have followed. The first
saving is multiplicative in the data, the second is additive.

\subsection{Longer Paths Do Not Pay}

For a path, the excluded case can be quantified exactly. Write $D = B(A)$ for the
delegated part and $P$ for the mappings over a smallest connected cover that satisfy the
anchor conditions; each $p \in P$ extends to one or more members of $D$, so
$|P| \leq |D|$.

\begin{corollary}[Joins removed]
Delegating a block removes at most $|B| - \gamma_c(B)$ joins per delegated mapping, and
attains this only when the plan forms a smallest connected cover. For a star of $k$
patterns the bound is $k - 1$, the whole block; for a path of length $n \geq 2$ it is $1$,
independent of $n$.
\end{corollary}

The bound is not always attained, because an adaptive plan chooses its own join order. Of
the two connected sub-blocks of size $n-1$ in a path --- the prefix $tp_1 \dots tp_{n-1}$
and the suffix $tp_2 \dots tp_n$ --- only the prefix contains $?s$ and hence covers the
anchors. An order that builds the suffix first reaches the anchor cover only at the root,
where no join is left to skip and pruning coincides with deduplication.

\begin{proposition}[Cost of delegating a path of length $n \geq 3$]
Let $c_{\mathit{req}}$ be the cost of a request, $c_{\mathit{parse}}$ the cost of
admitting one mapping from the resource, $c_{\mathit{join}}$ one join probe and
$c_{\mathit{mat}}$ one materialised mapping. Then
\[
W_{CR} - W_{\mathit{base}} \;=\;
c_{\mathit{req}} + |D|\,c_{\mathit{parse}}
\;-\; \bigl( |P|\,c_{\mathit{join}} + |D|\,c_{\mathit{mat}} \bigr),
\]
where the subtracted term is present only for the mappings whose join order forms a
smallest connected cover, and is zero otherwise.
\end{proposition}

\begin{proof}
Everything downstream of $B$ is unchanged: $|D|$ mappings enter the rest of the query
whether the resource or the default plan produced them. Within $B$, the default plan
performs the same joins as without the resource up to depth $\gamma_c(B) = n-1$, since no
shallower operator can prune. For the mappings that reach a cover, the final join and the
materialisation of their completions are skipped; for the others nothing is skipped. The
resource adds one request and the admission of $|D|$ mappings.
\end{proof}

Since $c_{\mathit{parse}} \geq c_{\mathit{mat}}$ --- admitting a mapping off the wire is
not cheaper than producing one from a hash probe --- the difference is non-negative in
every join order. Delegating a path of length three or more therefore never reduces total
work. What it reduces is latency: the default plan needs $n$ dependent round trips, since
the document binding $?o_{i}$ cannot be fetched before $?o_{i-1}$ resolves, whereas the
resource answers in one. The additional work is performed while the engine would
otherwise be blocked, so it is free in wall-clock terms exactly while the engine is
I/O-bound, and not otherwise.

\subsection{Reducible Decompositions}

Excluding long paths from delegation does not mean answering them without composite
resources, because a path can be partitioned into reducible blocks.

\begin{proposition}[Tiling]
Every connected sub-block of a path $tp_1, \dots, tp_n$ is a contiguous run, and a run of
$m$ patterns is reducible if and only if $m \leq 2$. Consequently every path admits a
partition into reducible blocks, and the smallest such partition has $\lceil n/2 \rceil$
blocks.
\end{proposition}

\begin{proof}
A run $tp_i, \dots, tp_j$ of length $m = j - i + 1$ has anchors
$\{o_{i-1}, \dots, o_{j-1}\}$ and $\gamma_c = m - 1$ for $m \geq 2$ by the earlier
proposition, and $\gamma_c = 1$ for $m = 1$; so it is reducible exactly when $m \leq 2$.
A partition into reducible blocks therefore uses blocks of at most two patterns and needs
at least $\lceil n/2 \rceil$ of them, and the partition into consecutive pairs, with one
singleton when $n$ is odd, attains that bound.
\end{proof}

The tiles are pattern-disjoint, so no two of them claim the same join entry, and each is
delegated under its own anchor conditions --- possibly to different pods, when the path
crosses pod boundaries. Their queries are parameterised rather than bound, so they do not
depend on one another and are requested concurrently: latency remains one round trip, as
it would be for a single resource covering the whole path, while every tile prunes at
depth $1$ and so removes its part of the delegated data from the default plan's join
state.

The tiles are, however, less selective than the path they decompose. A resource answering
$tp_1 \dots tp_4$ returns the mappings of the whole path, whereas two tiles return the
mappings of $tp_1 tp_2$ and of $tp_3 tp_4$ separately and leave their join to the client.
Decomposition therefore trades transferred volume and client-side join work for input
pruning and lower per-request server compute. Which side wins depends on the selectivity
of the path relative to its halves, and is an empirical question rather than one the
analysis settles.

\subsection{Guidance for Publishers}

Because reducibility is a property of the published block rather than of the client, the
criterion also says what a pod should expose. A resource answering a long path gives
clients a latency improvement and no reduction in work, since its result is rebuilt by
the client's default operators before being discarded. The same data published as
reducible tiles gives both, at the cost of one request per tile and a less selective
answer. A pod that publishes composite resources for its data is therefore better served
by publishing blocks of unit cover number --- stars over its subjects, and two-hop paths
--- than by publishing the longest paths it can compute.

\subsection{What the Analysis Predicts}

The account above is falsifiable on three counts, and an evaluation that delegates long
paths as an ablation can check each directly:

\begin{enumerate}
\item time to first result improves for delegated paths of any length, in proportion to
      the number of dependent round trips removed;
\item total work increases for paths of length three or more, and decreases for stars and
      two-hop paths;
\item the distribution of pruning depth concentrates at depth $1$ for reducible blocks and
      at or near the root for longer paths, with the fraction reaching depth $n-1$
      governed by how often the adaptive order builds the prefix rather than the suffix.
\end{enumerate}

The third is the most informative, as it separates the two mechanisms: a long path whose
prunes all land at the root is contributing nothing but a request and a parse, and the
measurement says so without reference to the timings.
